% ============================================================
% CHAPITRE 3 : IMPLÉMENTATION
% ============================================================
\chapter{Implémentation}
\label{ch:impl}

\section{Architecture logicielle}

Le code source d'\AZIZA est structuré en modules indépendants dans une arborescence claire :

\begin{lstlisting}[caption=Structure du projet AZIZA]
AZIZA/
├── aziza/
│   ├── __init__.py
│   ├── pipeline/
│   │   ├── complet.py       # Orchestrateur : fusion forme + fond + verrous
│   │   ├── forme.py         # Evaluation de la forme (/7)
│   │   ├── fond.py          # Evaluation du fond (/3) via Mistral
│   │   ├── similarite.py    # Metriques lexicales (LCS, parabolique)
│   │   └── seuils_production.json
│   └── machine_etats.py     # Machine a etats adaptative (3 paliers)
├── demo/
│   ├── app.py               # Routes HTML + Swagger + rate limiting
│   ├── api.py               # API REST (/api/v1) + auth
│   ├── static/style.css
│   ├── templates/           # 4 templates Jinja2
│   └── data/
├── docs/                    # Memoire LaTeX
├── experiments/             # 7 notebooks Jupyter + figures
├── tests.py                 # 117 tests unitaires + integration
├── requirements.txt
├── .env.example
└── README.md
\end{lstlisting}

\section{Implémentation du pipeline Forme}

La fonction \texttt{analyser\_forme} dans \texttt{forme.py} coordonne trois sous-systèmes :

\subsection{Lexical}

Le calcul du score lexical utilise le LCS normalisé :
\begin{equation}
\text{LCS\_norm} = \frac{\text{LCS}(source, contraction)}{\max(|source|, |contraction|)}
\end{equation}
Transformé par une fonction parabolique :
\begin{equation}
f(x) = -4.15x^2 + 6.64x - 1.66
\end{equation}
Bornée entre \(S_{\min}=0.26\) (score 0) et \(S_{\max}=0.80\) (score 3).

\subsection{Syntaxe et connecteurs}

Neuf catégories de connecteurs sont détectées par expressions régulières et comparées à un seuil de 3 connecteurs attendus dans une contraction de qualité.

\subsection{Orthographe}

L'appel à l'API LanguageTool distant renvoie la liste des erreurs avec leurs suggestions. Le barème dégressif pénalise proportionnellement au nombre de fautes.

\section{Implémentation du pipeline Fond}

\subsection{Défi technique : le format de l'API Mistral}

Lors du développement, un bug critique est apparu : l'API Mistral a modifié son format de message (2026). Le champ \texttt{content} doit désormais être un tableau de parties :

\begin{lstlisting}[caption=Correctif du format content pour Mistral API 2026,language=Python]
def _as_text(content: str) -> list:
    return [{"type": "text", "text": content}]
\end{lstlisting}

Ce correctif, appliqué dans \texttt{fond.py}, a résolu l'erreur \texttt{Cannot read "image.png"} qui survenait lorsque Mistral recevait une chaîne simple au lieu du format tableau.

\subsection{Embedding sémantique}

Le calcul de similarité cosinus utilise \texttt{mistral-embed} :

\begin{equation}
\text{sim\_sem} = \frac{\vec{v}_{source} \cdot \vec{v}_{contraction}}{\|\vec{v}_{source}\| \times \|\vec{v}_{contraction}\|}
\end{equation}

\subsection{Prompt Few-Shot avec chaîne de pensée}

Le prompt système impose une chaîne de pensée obligatoire en 6 étapes avant la production du JSON de sortie. Cinq exemples illustrent les cas typiques (correct, contresens, subjectivité, omissions).

\section{API REST et sécurité}

L'API REST, définie dans \texttt{demo/api.py}, expose deux endpoints :

\begin{itemize}
    \item \texttt{POST /api/v1/evaluer} : évaluation complète nécessitant le header \texttt{X-API-Key}.
    \item \texttt{GET /api/v1/health} : healthcheck public.
\end{itemize}

Le \textbf{rate limiting} (60 requêtes/min) est implémenté via \texttt{flask-limiter}, avec une clé de limitation basée sur \texttt{X-API-Key} ou l'adresse IP.

\section{Documentation Swagger}

La documentation interactive est générée automatiquement par \texttt{flasgger} à partir des docstrings YAML des routes. Un correctif a été nécessaire : le paramètre \texttt{headers} doit être défini à \texttt{[]} dans la configuration Swagger pour éviter un \texttt{TypeError} (`NoneType` object is not iterable).

\section{Déploiement}

Le déploiement initial prévu sur HuggingFace Spaces (Docker) a été abandonné suite au passage de Docker Spaces derrière un paywall en juillet 2026. L'application a été redéployée sur \textbf{PythonAnywhere} (gratuit, pas de sleep), accessible à l'adresse :
\begin{center}
\url{https://Jiovany11.pythonanywhere.com}
\end{center}

La procédure de déploiement inclut :
\begin{enumerate}
    \item Clone du dépôt GitHub.
    \item Création d'un environnement virtuel Python 3.12.
    \item Installation des dépendances (\texttt{requirements.txt}).
    \item Configuration manuelle du Web App (WSGI, virtualenv, static files).
    \item Définition des variables d'environnement (\texttt{MISTRAL\_API\_KEY}, \texttt{AZIZA\_API\_KEY}, \texttt{FLASK\_SECRET\_KEY}).
\end{enumerate}
